% UBT-AI-PROVENANCE-BEGIN
% schema: ubt-ai-provenance/v1
% tier: C_working
% ai_assistance: disclosed
% human_review: risk-based
% editorial_responsibility: Ing. David Jaroš
% policy: ../../../AI_PROVENANCE.md
% notice: Working material; exhaustive human review is not claimed.
% UBT-AI-PROVENANCE-END

\chapter{UBT za hodinu: motivace, architektura a současný stav}
\label{ch:overview}

\section{K čemu tato kapitola slouží}
Unified Biquaternion Theory (UBT) se ptá, zda geometrie prostoročasu a vnitřní
Struktura pole může být zakódována do jednoho komplexního hlavního pole s kvaternionovou hodnotou,
spíše než představeny jako nesouvisející základní objekty.  Aktivní spouštění
zápis je
\[
 \Theta(q,\tau)\in\mathbb B=\mathbb C\otimes\mathbb H,
 \qquad \tau=t+i\psi .
\]
Zápis je kompaktní, ale skrývá několik logicky odlišných otázek:
co je definováno, co je kinematicky reprezentovatelné, co vyplývá z an
akce a co je pouze výzkumná hypotéza.  Učebnice je uchovává
úrovně jsou v celém rozsahu odděleny.

\section{Současný lokální geometrický řetězec}
Čistá místní větev GR začíná od kovariančního derivátu singlu
pole,
\[
 E_\mu:=\mathcal N_0^{-1/2}D_\mu\Theta .
\]
Napište na skutečný Lorentzův plátek
\[
 E_\mu=i e_\mu{}^0\mathbf1+e_\mu{}^k\mathbf e_k,
 \qquad e_\mu{}^a\in\mathbb R.
\]
Kvaternionová konjugace je označena $\sharp$.  Symetrický produkt je
centrální a definuje metriku,
\[
 \boxed{
 \frac12\left(E_\mu^\sharp E_\nu+E_\nu^\sharp E_\mu\right)
 =g_{\mu\nu}\mathbf1 .}
\]
Minimální místní metrická konstrukce tedy nevyžaduje pozorovatele
projekce, průměr z kompaktních vláken nebo mapa pro vkládání.  Toto nahrazuje an
starší architektonický snímek dochovaný v historickém archivu.

Stejná první derivace se chová jako tetráda.  Pro nedegenerovanou tetrádu
mapa $e_\mu{}^a\mapsto g_{\mu\nu}$ má známé počítání
\[
 16-6=10,
\]
protože šest lokálních Lorentzových změn ponechává deset metrických komponent neměnných.
To řeší problém místního pořadí na kinematické úrovni.

\section{Konexe, křivost a co je stále dynamické}
Místní Lorentz/bikvaternionické spojení vstupuje do $D_\mu$.  V nekrouceném GR
Odbočte kompatibilní spinové připojení je Levi--Civita spinové připojení
počítáno z tetrády.  S kroucením,
\[
 \omega=\mathring\omega(e)+K(T),
\]
takže zadání tetrády a torze určuje metricky kompatibilní
spojení.  Zakřivení patří ke komutátoru kovariančních derivátů,
ne na algebraický komutátor nohou tetrády:
\[
 [D_\mu,D_\nu]\quad\hbox{contains curvature.}
\]

Současná GR práce uzavírá důležité otázky místní reprezentativnosti, ale je
ještě neodvozuje veškerou Einsteinovu dynamiku z původní mistrovské akce UBT.
Zejména výběr akce, omezení fyzického kroucení, plný
poruchový můstek a globální pokračování zůstávají stavově řízené otevřené
problémy.  Směrodatné znění je zachováno v
\texttt{STATUS.md} a \texttt{WHAT\_IS\_PROVED.md}.

\section{Komplexní čas a theta/spektrální větev}
Kanonický zápis zůstává zachován
\[
 \tau=t+i\psi,
\]
s $\psi$ interpretovaným v aktivní komplexní časové větvi jako vnitřní fáze
spíše než druhé makroskopické hodiny.  Pak zmenšené theta modely
přirozeně obsahují faktory formy
\[
 S(t,\psi)=\sum_n a_n e^{-\pi\psi n^2}e^{i\pi t n^2}.
\]
Tento jeden výraz připouští několik standardních reprezentací: Fourier v $t$,
Laplace nebo Mellin v $\psi$, Jacobi/Poisson dualitě a v řízeném
kvadratické příklady, vztah mezi součty režimů a součty dráhy/vinutí.  Tito
klasické nástroje jsou vyvinuty výslovně později, protože organizují několik
jinak odpojené výzkumné dráhy UBT.

\section{Kalibrační a částicové sektory}
Algebru $\mathbb C\otimes\mathbb H\cong M_2(\mathbb C)$ přirozeně podporuje
levé a pravé akce, spinorální Lorentzovy transformace a vnitřní
symetrické konstrukce.  Některá měřidla-sektorová tvrzení jsou prokázána popř
podmíněné na uvedených úrovních; ostatní, včetně částí hypercharge,
výběr generací a dynamické narušení symetrie zůstávají otevřené.  Tento
učebnice proto nikdy neodvozuje úplné odvození standardního modelu pouze z
existence biquaternionové algebry.

\section{Konstanta jemné struktury: současné bezpečné tvrzení}
Program alfa obsahuje reprodukovatelný strukturální mechanismus, který vybírá
celé číslo $137$ za uvedených předpokladů, včetně výsledku prvočísla stability.
Celková trať však v současné době není odvozením prvních principů
fyzikální konstanta jemné struktury.  Zvláště důležité jsou dva blokátory:
\begin{enumerate}
 \item efektivní koeficient požadovaný potenciálem vinutí nebyl splněn
       odvozeno jedinečně z působení UBT (Gap G137-B);
 \item multiplicita auditované smyčky a multiplicita twist-sektoru nejsou
       přesto identifikovaný prokázaným mostem.
\end{enumerate}
V důsledku toho ani fráze ``odvození alfa bez fit'' ani numerické
hodnota $137.036\ldots$ by zde měla být prezentována jako uzavřená predikce UBT.
Kapitola~\ref{ch:alpha} odvozuje podmíněnou matematiku a uvádí mezery
na každém kroku.

\section{Jak číst statusové značky}
Úložiště používá dva nezávislé účetní systémy, které nesmí být
zmatený.

Za prvé, úrovně vědeckých tvrzení rozlišují standardní identity, prokázané UBT
lemmata, podmíněné výsledky, motivované dohady, pozorování a otevřené
mezery.  Za druhé, úrovně původu AI zaznamenávají stav redakce/revize souborů.
Ty jsou definovány \texttt{PROVENANCE\_TIERS.yaml}:
Úroveň A je ověřena autory, úroveň B strojově ověřená, úroveň C funkční a úroveň D
historický.  Tato učebnice je pracovním materiálem úrovně C.  Jeho úlohou je vysvětlovat
a propojit matematiku, ne tiše prosazovat nárok na úroveň A.

\section{Pravidlo autoritativního zdroje pro učebnici}
Zdrojová mapa v \texttt{docs/textbook/SOURCE\_MAP.md} zaznamenává, kde každý
kapitola získá svůj živý stav.  Historický materiál v \texttt{ARCHIV/} může být
cenné pro pochopení toho, jak se nápad vyvíjel, ale nikdy není tichý
\texttt{\textbackslash input} jako součást aktuální teorie.
